% third23 paper draft — THE DEFECT-3 PAPER (opened 2026-08-16; matured 2026-08-18 after
% round 38 banked the Three-Defect Chain Theorem). Source of truth for every theorem: the
% campaign wiki entries + ASSEMBLY_LEDGER.md named in comments; every proof here is
% assembled from those banked artifacts, never re-derived. Certificate corpora are cited
% as machine-checkable proof objects with SHA-256 anchors (Section on proof objects).
% Gates before submission: principal line-read + external cold-read referee.
% TODO(acquire): Aigner 1985 deposited copy for the width-two equality citation.
\documentclass[11pt]{article}
\usepackage{amsmath,amssymb,amsthm}
\usepackage{booktabs}
\usepackage[margin=1.2in]{geometry}

\newtheorem{theorem}{Theorem}[section]
\newtheorem{lemma}[theorem]{Lemma}
\newtheorem{proposition}[theorem]{Proposition}
\newtheorem{corollary}[theorem]{Corollary}
\theoremstyle{definition}
\newtheorem{definition}[theorem]{Definition}
\theoremstyle{remark}
\newtheorem{remark}[theorem]{Remark}

\newcommand{\pr}{\operatorname{Pr}}
\DeclareMathOperator{\bal}{bal}

\title{The $1/3$--$2/3$ Conjecture on the Defect Ladder:\\
The Two- and Three-Defect Chain Theorems}
\author{}% TODO(principal): authorship
\date{Draft of August 18, 2026}

\begin{document}
\maketitle

\begin{abstract}
A finite poset is a \emph{$k$-defect chain} if it becomes a chain after the deletion of
$k$ elements; equivalently, its maximum chain omits exactly $k$ elements. We prove the
$1/3$--$2/3$ conjecture for all posets of defect at most three: every such poset that is
not itself a chain has an incomparable pair $(u,v)$ with $1/3\le\pr(u\prec v)\le 2/3$.
The class contains posets of width three and four, does not bound the number of
incomparabilities per element, and we exhibit members escaping each previously proved
class known to us. The two-defect case is proved by an explicit balance lemma for
overlapping insertion windows; for it we also prove equality rigidity: every
ordinal-sum-indecomposable poset of defect exactly two has a strictly balanced eligible
count, and among chain-plus-two-incomparable-defect presentations the unique non-strict
member is the three-element poset $Q$ (itself of defect one) --- extending Aigner's
width-two equality classification to an infinite family of width-three posets. The
three-defect case is proved by contradiction through a rigidity theory of \emph{closed-third
avoidance}: if every incomparable orientation probability avoided the closed interval
$[N/3,2N/3]$, the majority relation would be a strict total order, every defect would
carry a unique \emph{heavy gap} holding more than a third of the extension mass, and the
heavy gaps would align with the defect order. An exhaustive partition of the surviving
configurations by defect shape (antichain, $V$, $\Lambda$, $Q$) and heavy-gap pattern is
then closed cell by cell: by Helly-type coincidence and marginal monotonicity arguments,
by weight-preserving reflection dualities, by vacuity theorems, and --- for the hardest
cells --- by exhaustive weak-order chamber decompositions carrying exact rational
nonnegative-coefficient certificates ($169$ chambers for the antichain common-heavy
inequality; $31$ strict and $31$ boundary chambers with $62$ signed certificates for the
tied-siblings $V$ cell and its one-relation corrections). All certificate corpora are
exact, machine-checkable, independently replayed proof objects anchored by cryptographic
hashes.
\end{abstract}

\section{Introduction}
% Positioning per the banked novelty sweep (wiki: operator-novelty-sweep-two-defect-
% theorem-new-rigidity-width). All class attributions verified against produced sources.
The $1/3$--$2/3$ conjecture --- posed independently by Kislitsyn \cite{kislitsyn},
Fredman \cite{fredman}, and Linial \cite{linial} in connection with sorting under partial
information --- asserts that every finite poset $P$ that is not a chain contains an
incomparable pair $(u,v)$ whose orientation probability under the uniform distribution on
linear extensions satisfies $1/3\le\pr(u\prec v)\le 2/3$. Kahn and Saks \cite{kahn-saks}
proved the first uniform positive bound; the strongest general bound is due to Brightwell,
Felsner, and Trotter \cite{bft}. The conjecture is known for partial orders of width two
(Linial \cite{linial}, sharpened by Sah \cite{sah}), height two \cite{tgf}, posets with at
most $14$ elements \cite{deloof,gupta}, posets in which every element is incomparable to
at most six others (Peczarski \cite{peczarski-thin}), series-parallel posets, $N$-free
posets and polytrees (Zaguia \cite{zaguia-nfree,zaguia-polytree}), semiorders (Brightwell
\cite{brightwell-semi}), several lattice and Young-diagram families (Olson and Sagan
\cite{olson-sagan}), and is preserved by lexicographic sums \cite{dolores-cuenca}. Chan
and Pak \cite{chan-pak} survey the area.

Every finite poset becomes a chain after deleting some number of elements; call the least
such number its \emph{defect}. Since the retained chain of a minimal deletion is a
maximum chain, the defect of $P$ equals $|P\setminus C|$ for any maximum chain $C$. The
conjecture restricted to defect $k$ interpolates, as $k$ grows, from the merging problems
of the classical literature to the full conjecture. Defect one is contained in width two.
This paper settles defects two and three completely, and classifies the extremal cases at
defect two.

Three subfamilies of the two-defect class were previously known and are not claimed here:
members whose two defect windows coincide carry a defect-exchanging automorphism, a case
settled by Peczarski's automorphism theorem \cite{peczarski-aut} (and, for cycle-free
members, already by Ganter, H\"afner, and Poguntke \cite{ghp}, who obtain a $1/2$-balanced
pair); members of width two are Linial--Sah; members whose windows have
length at most two are semiorders. The class as a whole --- with windows of unbounded length, hence unboundedly many
incomparabilities at a single element, and with genuine width-three members --- is
contained in neither the width-two class nor any bounded-incomparability class, and its
non-containment in the remaining listed classes is witnessed explicitly. With $m=2$ and
windows $I=[0,1]$, $J=[1,2]$, the four elements $(x,c_1,c_2,y)$ induce an $N$ ($x<c_2$,
$c_1<c_2$, $c_1<y$, all other pairs incomparable), both as an induced suborder and in
the Hasse diagram --- so the class is contained in neither the series-parallel nor the
$N$-free class. With $m=3$ and windows $I=[1,2]$, $J=[1,3]$, the Hasse diagram contains
the cycle $c_1\,x\,c_3\,c_2$ --- so it is not contained in the forest-cover class.
Unbounded windows exclude semiorders, and height excludes the height-two class. The incomparable-defect
members are precisely the interval orders with exactly two nontrivial intervals; general
interval orders remain open.

Section~2 constructs the insertion-window model and its state bijection; Section~3
proves the balance lemma; Section~4 assembles the Two-Defect Chain Theorem; Section~5
proves the equality rigidity classification; Sections~6--7 enter the three-defect ladder
through its exactly solvable chambers (the aligned $Q$ chamber, exactly at the boundary,
and the adjacent one-endpoint chambers). Sections~8--13 prove the Three-Defect Chain
Theorem: Section~8 develops the closed-third avoidance rigidity theory (strong order,
heavy gaps, alignment), Section~9 the weighted gap-fiber model and the structural
reductions, Sections~10--12 close the twelve shape-by-heavy-pattern cells (antichain;
$V$ and $\Lambda$; $Q$), and Section~13 assembles the theorem. Section~14 documents the
certificate corpora as verifiable proof objects, and Section~15 discusses the ladder
ahead. Throughout, $e(P)$ is the number of linear extensions of $P$, $\pr(u\prec v)$ the
fraction orienting $u$ before $v$, and $\#L(P;\,\cdot)$ the number of extensions with the
displayed property; windows are integer intervals of insertion gaps against the fixed
chain $C$.

\section{The doubled-diagonal model}
% Source: wiki doubled-diagonal-overlapping-interval-balance-decision (round 5, sol),
% independently re-derived in round 6 (supervisor-verdict-work-d4e5c56a7791).
Let $C=(c_1<\cdots<c_m)$ be a chain and let $x,y$ be two further elements with
$x\parallel y$. Consistency with the chain forces the admissible insertion gaps of $x$ to
form a window $I=[a,b]\subseteq\{0,\dots,m\}$ (gap $i$ meaning that exactly $i$ chain
elements precede $x$): the relations are $c_t<x$ for $t\le a$ and $x<c_{t}$ for $t>b$.
Similarly $y$ carries a window $J=[c,d]$.

\begin{lemma}[Overlap]\label{lem:overlap}
If $x\parallel y$ then $I\cap J\neq\emptyset$. Conversely, for any two nonempty
overlapping windows $I,J$, imposing the chain order, the threshold relations of both
windows, and no relation between $x$ and $y$ defines a valid poset in which $x\parallel y$
and no hidden comparabilities arise.
\end{lemma}

\begin{proof}
If $b<c$ then $x<c_{b+1}\le c_c<y$ by transitivity, contradicting $x\parallel y$; the case
$d<a$ is symmetric. For the converse, every pair of gaps $(i,j)\in I\times J$ is realized
by merging $x$ and $y$ into the chain at those gaps, and at a common gap both relative
orders of $x,y$ occur, so $x\parallel y$ and the prescribed windows are exactly attained.
% imported from: round-6 verdict items "Insertion-Gap Structural Model" and "Overlap
% Converse for Two Insertion Intervals"
\end{proof}

\begin{lemma}[State bijection and counts]\label{lem:states}
Let $p=|I|$, $q=|J|$, $K=I\cap J$, $h=|K|$. The linear extensions of $P$ are in bijection
with the states $(i,j)\in I\times J$ with $i\ne j$, together with two ordered states for
each common gap $i=j\in K$; hence
\[
N=e(P)=(pq-h)+2h=pq+h.
\]
The orientation counts of the incomparable pairs are
\begin{align*}
D&=\#\{(i,j)\in I\times J: i<j\}+h &&\text{for the pair }(x,y),\\
A_t&=q\,\#\{i\in I: i<t\}+\#\{i\in K: i<t\} &&\text{for }(x,c_t),\ a<t\le b,\\
B_t&=p\,\#\{j\in J: j<t\}+\#\{j\in K: j<t\} &&\text{for }(y,c_t),\ c<t\le d,
\end{align*}
and exactly the displayed thresholds correspond to incomparable pairs; outside them the
orientation is forced.
\end{lemma}

\begin{proof}
A linear extension is determined by the pair of gaps, except that when $x$ and $y$ occupy
a common gap the two relative orders give two extensions --- the doubled diagonal. For
$D$: every off-diagonal state with $i<j$ contributes, and exactly the $x$-before-$y$
diagonal state contributes, giving the displayed count. For $A_t$ with $a<t\le b$: $x<c_t$
exactly when the $x$-gap satisfies $i<t$; each such $i$ pairs with $q$ choices of $j$,
plus one extra state when $i\in K$ (the doubled diagonal), and symmetrically for $B_t$.
Incomparability of the displayed pairs is witnessed by explicit gap choices realizing both
orientations.
% imported from: r5 artifact section 1 (state model), verified against the m<=4
% topological-sort audit (259 cases, 0 failures) recorded in the same artifact
\end{proof}

\section{The balance lemma}
% Source: same r5/r6 entries; full derivation banked verbatim in the r5 artifact.
\begin{theorem}[Overlapping Two-Interval Balance Lemma]\label{thm:balance}
For every pair of nonempty overlapping windows, at least one of $D$, an eligible $A_t$,
or an eligible $B_t$ lies in $[N/3,\,2N/3]$.
\end{theorem}

The proof rests on one discrete intermediate-value fact and exact jump accounting.

\begin{lemma}[Discrete crossing lemma]\label{lem:crossing}
Let $x_1<\cdots<x_r$ be integers with $x_1\le N/3$, $x_r\ge 2N/3$, and
$x_{k+1}-x_k\le N/3$ for all $k$. Then some $x_k\in[N/3,\,2N/3]$.
\end{lemma}

\begin{proof}
Take the first $k$ with $x_k\ge N/3$. If $k=1$ then $x_1=N/3$ exactly and we are done;
otherwise $x_{k-1}<N/3$, so $x_k=x_{k-1}+(x_k-x_{k-1})<N/3+N/3=2N/3$.
\end{proof}

Exchanging the defect names replaces $D$ by $N-D$ and swaps the $A$ and $B$ families;
translation and order reversal also preserve the assertion. After these symmetries exactly
two canonical forms remain: \emph{crossing} ($a\le c\le b\le d$; then $h=p-s$ where $s=c$
after translating $a$ to $0$, and $q\ge h$) and \emph{nested} ($J\subseteq I$; then
$h=q$).

\paragraph{The crossing form.} Direct summation gives
\[
N=pq+h,\qquad D=pq-\tfrac{h(h-1)}{2},\qquad
A_t=qt+\max(0,t-s),\qquad B_{s+u}=pu+\min(u,h),
\]
with exact consecutive jumps $A_{t+1}-A_t\in\{q,q+1\}$ and
$B_{s+u+1}-B_{s+u}\in\{p,p+1\}$. If $N\ge 3(q+1)$ then $p\ge2$, the $A$-sequence runs
from $A_1\le q+1\le N/3$ to $A_{p-1}=N-(q+1)\ge 2N/3$ with jumps at most $q+1\le N/3$,
and Lemma~\ref{lem:crossing} produces a balanced eligible $A_t$. Symmetrically, if
$N\ge3(p+1)$ the $B$-sequence does the same. If neither large-case condition holds, the
inequalities $q(p-3)<3-h$ and $p(q-3)<3-h$ with $1\le h\le\min(p,q)$ force exactly fifteen
triples $(p,q,h)$, each of which carries an explicit balanced witness:

\begin{center}
\begin{tabular}{@{}lrrl@{}}\toprule
$(p,q,h)$ & $s$ & $N$ & witness \\ \midrule
$(1,1,1)$ & 0 & 2 & $D=1$\\
$(1,2,1)$ & 0 & 3 & $D=2$\\
$(1,3,1)$ & 0 & 4 & $B_{s+1}=2$\\
$(1,4,1)$ & 0 & 5 & $B_{s+1}=2$\\
$(2,1,1)$ & 1 & 3 & $D=2$\\
$(2,2,1)$ & 1 & 5 & $A_1=2$\\
$(2,2,2)$ & 0 & 6 & $D=3$\\
$(2,3,1)$ & 1 & 7 & $A_1=3$\\
$(2,3,2)$ & 0 & 8 & $D=5$\\
$(3,1,1)$ & 2 & 4 & $A_2=2$\\
$(3,2,1)$ & 2 & 7 & $A_2=4$\\
$(3,2,2)$ & 1 & 8 & $D=5$\\
$(3,3,1)$ & 2 & 10 & $A_2=6$\\
$(3,3,2)$ & 1 & 11 & $A_2=7$\\
$(4,1,1)$ & 3 & 5 & $A_2=2$\\ \bottomrule
\end{tabular}
\end{center}

\paragraph{The nested form.} Here $N=(p+1)q$. If $q\ge2$, take $u=\lceil q/3\rceil$: then
$1\le u\le q-1$ and $q\le 3u\le 2q$, so the eligible count $B_{s+u}=(p+1)u$ satisfies
$N\le 3B_{s+u}\le 2N$ \emph{exactly} --- the factorization of $N$ makes the nested case
rigid. If $q=1$ then $N=p+1$, $D=s+1$, and $A_t=t+\mathbf1_{\{s<t\}}$; for $p\ge5$ the
$A$-sequence starts at most $2$, ends at least $p-1$, and jumps by $1$ or $2$, so
Lemma~\ref{lem:crossing} applies, while $p\le4$ is a finite table of explicit witnesses.
% imported from: r5 artifact sections 4-5; the closed interval is used throughout and no
% continuity is invoked (r5 artifact section 6). Corroborating audits (r5/r6): all 5,551
% overlapping endpoint tuples at m=12; adversarial canonical scan of 4,040,100 states
% p,q<=200; zero failures. Audits are checks, not premises.

\begin{remark}[$D$ alone does not suffice]
The count $D$ alone is not always balanced: at $(p,q,h)=(2,2,1)$ one has $N=5$ and
$D=4>2N/3$, while the eligible $A_1=2$ is balanced. Any proof must use the chain-cut
counts. % source: r6 verdict, "D-Only Balance Route Refuted"
\end{remark}

\section{The two-defect chain theorem}
\begin{theorem}[Two-Defect Chain Theorem]\label{thm:twodefect}
Every finite non-chain poset $P$ containing a chain $C$ with $|P\setminus C|=2$ has an
incomparable pair $(u,v)$ with $1/3\le\pr(u\prec v)\le 2/3$.
\end{theorem}

\begin{proof}
Write $P\setminus C=\{x,y\}$. If $x$ and $y$ are comparable, then $C\cup\{x<y\}$ is a
two-chain cover, so $P$ has width at most two and the theorem of Linial \cite{linial}
(in the sharpened form of Sah \cite{sah}) applies. If $x\parallel y$, their windows
overlap by Lemma~\ref{lem:overlap}; Lemma~\ref{lem:states} identifies the orientation
counts of the genuine incomparable pairs, and Theorem~\ref{thm:balance} produces one in
the middle third.
\end{proof}

\section{The extremal classification}
% Sources: r7 (supervisor verdict "DD-EQ Rigidity Proved") + r8 independent re-derivation
% (supervisor-verdict-work-5fbb947628b0). ATTRIBUTION: the width-two part is Aigner
% \cite{aigner}: width-two posets with balance constant exactly 1/3 are the ordinal sums
% of singletons and copies of Q. The width-three content is new; the order-14 census
% \cite{gupta} observes it empirically and notes it is beyond Aigner's argument.
Write $Q$ for the three-element poset with exactly one relation. Aigner \cite{aigner}
proved that the width-two posets with balance constant exactly $1/3$ are precisely the
ordinal sums of singletons and copies of $Q$. The census of Gupta \cite{gupta}
determines the extremal balance data through order $14$: exactly $128$ classes attain
$1/3$, every one an ordinal sum of singletons and copies of $Q$ --- in particular all of
width two. The following theorem proves this non-attainment for the full (infinite,
width-three-containing) class. Note the distinction between the \emph{presentation}
(a chain plus two designated defects, the model of Section~2, which any poset of defect
at most two admits) and the \emph{exact defect} $n-h(P)$ of the abstract: the extremal
member $Q$ arises as a one-element chain plus two defects but has exact defect one.

\begin{theorem}[Equality rigidity]\label{thm:rigidity}
\emph{(a)} Every ordinal-sum-indecomposable poset of defect exactly two has an eligible
orientation count strictly inside the open interval $(N/3,\,2N/3)$; in particular its
balance constant exceeds $1/3$.
\emph{(b)} Among ordinal-sum-indecomposable posets presented as a chain plus two
incomparable defects, the unique member all of whose eligible counts avoid the open
interval is $Q$ --- which has defect one. Every other member has an eligible count
strictly inside.
\end{theorem}

The proof classifies the \emph{simultaneous avoiders} --- parameter tuples for which every
eligible count $C$ satisfies $3C\le N$ or $3C\ge2N$ --- and intersects them with an
indecomposability criterion:

\begin{proposition}[Avoider classification]\label{prop:avoiders}
In the canonical crossing form, simultaneous avoidance holds exactly for
$(p,q,h,s)=(1,2,1,0)$ and $(2,1,1,1)$. In the nested form, exactly for
$(p,q,s)=(2,1,0)$ and $(2,1,1)$.
\end{proposition}

\begin{proposition}[Indecomposability criterion]\label{prop:indec}
The two-defect model with windows $I=[a,b]$, $J=[c,d]$ is ordinal-sum-indecomposable iff
its incomparability graph is connected, iff every $t\in\{1,\dots,m\}$ satisfies $a<t\le b$
or $c<t\le d$; under the overlap hypothesis this is equivalent to $\min(a,c)=0$ and
$\max(b,d)=m$.
\end{proposition}

\begin{proof}[Proof of Proposition~\ref{prop:avoiders}]
% transcribed from paper/derivations/r7_capture_02.md (r7, sol), independently
% re-derived in r8 (r8_capture_01/02, luna) — two derivations on record.
Write $\mathrm{AV}(C)$ for the condition $3C\le N$ or $3C\ge2N$. We first sharpen
Lemma~\ref{lem:crossing} under avoidance: if $x_1\le N/3$, $x_r\ge2N/3$, all jumps are at
most $N/3$, and \emph{every} term satisfies $\mathrm{AV}$, then the first term
$\ge N/3$ equals $N/3$ exactly, and its successor equals $2N/3$ exactly (the predecessor
of the first crossing term is strictly below $N/3$, so the jump bound puts that term
strictly below $2N/3$; avoidance then forces equality at $N/3$, and the next term can
neither be in the open middle nor overshoot $2N/3$).

\emph{Large regimes.} In the crossing form the $A$-sequence has $A_1\le q+1$,
$A_{p-1}=N-(q+1)$, and jumps in $\{q,q+1\}$. If $N\ge3(q+1)$, all-$A$ avoidance forces
$N/3\in\{q,q+1\}$ by the sharpened lemma, hence $N=3(q+1)$, i.e.\ $q(p-3)=3-h$. With
$1\le h\le\min(p,q)$ the solutions are $(p,q,h)=(3,q,3)$ for $q\ge3$, $(4,2,1)$, and
$(5,1,1)$; the first is killed at $q=3$ by $D=6$, $N=12$ (strict middle) and for $q>3$
by the symmetric $B$-argument (which would force $N=3(p+1)=12$, impossible), while the
last two carry the strict witnesses $A_2=4$, $N=9$ and $A_3=3$, $N=6$. Symmetrically,
$N\ge3(p+1)$ with all-$B$ avoidance forces $p(q-3)=3-h$, whose solutions
$(p,3,3)_{p\ge3}$, $(2,4,1)$, $(1,5,1)$ are killed the same way. Hence every survivor
satisfies $N<3(q+1)$ and $N<3(p+1)$.

\emph{Residual list.} The strict inequalities $q(p-3)<3-h$ and $p(q-3)<3-h$ force
$p,q\le4$ and exactly the fifteen triples of the table in
Section~\ref{thm:balance}. All but two carry an eligible count strictly inside the open
middle third:
\[
\begin{array}{llllll}
(1,1,1)\!: D{=}1,N{=}2 & (1,3,1)\!: B_1{=}2,N{=}4 & (1,4,1)\!: B_1{=}2,N{=}5\\
(2,2,1)\!: A_1{=}2,N{=}5 & (2,2,2)\!: D{=}3,N{=}6 & (2,3,1)\!: A_1{=}3,N{=}7\\
(2,3,2)\!: D{=}5,N{=}8 & (3,1,1)\!: A_2{=}2,N{=}4 & (3,2,1)\!: A_2{=}4,N{=}7\\
(3,2,2)\!: D{=}5,N{=}8 & (3,3,1)\!: A_2{=}6,N{=}10 & (3,3,2)\!: A_2{=}7,N{=}11\\
(4,1,1)\!: A_2{=}2,N{=}5
\end{array}
\]
(each displayed value $C$ satisfies $N<3C<2N$). The remaining two triples avoid: at
$(1,2,1)$, $N=3$ with $D=B_1=2$, both at the boundary $3C=2N$; at $(2,1,1)$, $N=3$ with
$D=2$ and $A_1=1$, at the boundaries $3C=2N$ and $3C=N$. Note $D-N/2=(pq-h^2)/2\ge0$
throughout, so $D$ never realizes the low branch in this orientation.

\emph{Nested form.} Here $N=(p+1)q$ and $B_{s+u}=(p+1)u$. If $q\ge2$ and $3\nmid q$,
then $u=\lceil q/3\rceil$ gives $q<3u<2q$, a strict-middle $B$; if $q=3k\ge6$, take
$u=k+1$. If $q=3$: $p=3$ gives $D=6$, $N=12$ strict middle, and for $p\ge4$ the
$A$-sequence starts at most $4<N/3=p+1$, ends at least $3p-1>2N/3$, with jumps at most
$4$, so Lemma~\ref{lem:crossing} gives a strict term. For $q=1$: $N=p+1$ and
$\{D\}\cup\{A_t\}=\{1,\dots,p\}$; $p=1$ fails at $D=1$, $N=2$, $p\ge3$ contains an
integer strictly inside, and $p=2$ survives for $s\in\{0,1\}$ with values $\{1,2\}$ at
the boundaries of $N=3$.
\end{proof}

\begin{proof}[Proof of Proposition~\ref{prop:indec}]
The incomparability graph has the edge $xy$, no edges among chain vertices, and $c_t$
adjacent to $x$ exactly for $a<t\le b$, to $y$ exactly for $c<t\le d$. After exchange
assume $a\le c$; overlap gives $c\le b$, so the union of the two adjacency windows is the
single index interval $(a,\max(b,d)]$, and the graph is connected exactly when that
interval is all of $\{1,\dots,m\}$ --- that is, $\min(a,c)=0$ and $\max(b,d)=m$.
Connectedness of the incomparability graph is equivalent to ordinal-sum
indecomposability, since the canonical ordinal-sum factors are precisely the
incomparability components. % banked: canonical-ordinal-sum-factorization theorem
\end{proof}

Intersecting the avoider tuples (with their translations, defect exchanges, and order
reversals) with the criterion leaves exactly the four interval presentations
$(L,F),(F,L),(F,R),(R,F)$, where $L=[0,0]$, $R=[1,1]$, $F=[0,1]$ at $m=1$: endpoint
coverage forces $m=s+q-1$, collapsing both survivor tuples to one-element chains with
the displayed windows, each isomorphic to $Q$ (extensions $abc,acb,cab$; pair counts
$(2,1)$ and $(1,2)$). Every other arithmetic survivor has a nontrivial chain prefix or
suffix and is a chain--$Q$--chain ordinal sum, hence decomposable. This proves
part~(b) of Theorem~\ref{thm:rigidity}.

Part (a) follows. Let $P$ be indecomposable of defect exactly two, and let $C$ be a
maximum chain, $P\setminus C=\{x,y\}$. If $x$ and $y$ are comparable, $C\cup\{x<y\}$ is
a two-chain cover, so $P$ has width two; Linial \cite{linial} gives $\delta(P)\ge1/3$,
and by Aigner's classification \cite{aigner} equality holds only for ordinal sums of
singletons and copies of $Q$ --- all decomposable except $Q$ itself, whose defect is
one. Hence $\delta(P)>1/3$: some incomparable pair lies strictly inside. If $x\parallel
y$, part (b) applies to the presentation, and its unique avoider $Q$ has defect one,
not two; so $P$ carries a strict eligible count.

\begin{remark}
An eligible count with $N/3<C<2N/3$ yields a strict witness; if all eligible counts avoid
the open interval, Theorem~\ref{thm:balance} forces equality at $N/3$ or $2N/3$ --- which
is the situation classified above. A corroborating exact scan through $p,q\le60$ (73{,}810
canonical cases) found no further avoiders. % r7; audits are checks, not premises
\end{remark}

\section{Three defects: the aligned chamber}
% Sources: r9 (supervisor verdict work_baac3d58cfc2), independently re-derived in r10
% (supervisor-verdict-work-98553ea81246; wiki aligned-q-defect-chain-theorem and
% siblings). DECIDED: two independent derivations.
We now take three defects with the induced order $Q$ itself: $x<y$, and $z$ incomparable
to both. Suppose first that all three insertion windows coincide in $I=[a,b]$ with
$p=|I|\ge2$.

\begin{lemma}[Triangular weighted-state model]\label{lem:tristate}
The linear extensions are represented exactly by the triples $(i,j,k)\in I^3$ with
$i\le j$ ($i,j,k$ the gaps of $x,y,z$), with multiplicity $1$ when the gaps are distinct
or $i=j\ne k$, multiplicity $2$ when $k$ ties exactly one of $i,j$, and multiplicity $3$
when $i=j=k$. In particular
\[
N=e(P)=\frac{p(p+1)(p+2)}{2}.
\]
\end{lemma}

\begin{theorem}[Aligned $Q$-Defect Chain Theorem]\label{thm:aligned}
In the aligned model, $\pr(x\prec z)=\pr(z\prec y)=2/3$ and
$\pr(z\prec x)=\pr(y\prec z)=1/3$ exactly.
\end{theorem}

\begin{proof}[Proof (orbit argument)]
Delete the relation $x<y$. The three defects become interchangeable, and relabeling
symmetry partitions the extensions of the weakened poset equally among the six relative
orders of $\{x,y,z\}$, each class of size $\binom{p+2}{3}$. Restoring $x<y$ retains
exactly the classes $xyz$, $xzy$, $zxy$. Two of the three have $x\prec z$, and two have
$z\prec y$; the counts are exact thirds of $N=3\binom{p+2}{3}$.
% imported from: r9 verdict "Aligned Q Orbit Enumeration", re-derived r10
\end{proof}

\begin{proposition}[Chain-cut counts]\label{prop:cuts}
For an eligible cut $c_t$ with $r=t-a\in\{1,\dots,p-1\}$, set
$H_1=\binom{p+2}{3}-\binom{p-r+2}{3}$, $H_2=r(r+1)(3p-2r+2)/6$, $H_3=\binom{r+2}{3}$.
Then $\#L(x<c_t)=2H_1+H_2$, $\#L(y<c_t)=H_2+2H_3$, $\#L(z<c_t)=H_1+H_2+H_3$, with
complements $N-{}$these; outside the eligible interval the orientation is forced.
\end{proposition}

\begin{remark}
The aligned family generalizes $Q$ ($p=2$, $m=1$ yields $Q$ together with one isolated
element --- the chain element $c_1$, incomparable to all three defects) and sits exactly on
the conjecture's boundary at its defect pairs for \emph{every} window size --- an infinite
three-defect family of boundary-exact pairs. Whether any member's balance \emph{constant}
equals $1/3$ is determined by Proposition~\ref{prop:cuts} and left to the discussion.
\end{remark}

\section{Three defects: the adjacent one-endpoint chambers}
% Source: r10 (wiki adjacent-one-endpoint-feasible-state-model / -extension-totals /
% -q-balance-theorem / -eligible-cut-atlas). Single-pair derivation with explicit
% witness algebra; verification status per the record.
Now break the symmetry minimally: keep the windows of $x,y$ equal to $I$ ($p\ge2$) and
let the window $J$ of $z$ be $I$ with one endpoint gap added or deleted.

\begin{lemma}[Feasible states]\label{lem:onepstate}
In every admissible adjacent one-endpoint chamber the linear extensions are exactly the
triples $(i,j,k)$ with $i,j\in I$, $i\le j$, $k\in J$, with weight
$w(i,j,k)=1+\mathbf1_{\{k=i\}}+\mathbf1_{\{k=j\}}$. With $M=\binom{p+2}{3}$,
$T=\binom{p+1}{2}$, $S=\binom{p+2}{2}$, the totals are
\[
N_+=3M+T=\frac{p(p+1)(p+3)}{2}\quad\text{(endpoint addition)},\qquad
N_-=3M-S=\frac{(p-1)(p+1)(p+2)}{2}\quad\text{(endpoint deletion, $p\ge2$)}.
\]
\end{lemma}

\begin{theorem}[Adjacent One-Endpoint $Q$ Balance Theorem]\label{thm:onep}
Every admissible adjacent one-endpoint chamber has a balanced incomparable defect pair.
Explicitly: right endpoint addition admits the witness $\{y,z\}$ with
$\pr = \frac{p+5}{3(p+3)}$; left endpoint addition admits $\{x,z\}$ with
$\pr = \frac{2(p+2)}{3(p+3)}$; right endpoint deletion admits $\{x,z\}$ with
$\pr = \frac{p(2p+5)}{3(p+1)(p+2)}$; left endpoint deletion admits $\{y,z\}$ with
$\pr = \frac{p^2+4p+6}{3(p+1)(p+2)}$. Each ratio lies in $[1/3,\,2/3]$ for every
$p\ge2$.
\end{theorem}

\begin{proof}
% transcribed from paper/derivations/r10_capture_03.md (the complete orientation table
% with polynomial numerators; chain-cut atlas in the sibling banked entry).
Write $R_x=\#L(z\prec x)$ and $R_y=\#L(z\prec y)$. Summing the weighted states of
Lemma~\ref{lem:onepstate} fiber by fiber (a tied gap $k=i$ contributes both tied orders
to $z\prec y$; the triple diagonal contributes its two $z$-before orders to $R_y$ and one
to $R_x$) gives
\[
R_x=\sum_{i=0}^{p-1}(p-i)\,\#\{k\in J:k<i\}+\sum_{i\in K}(p-i),\qquad
R_y=\sum_{j=0}^{p-1}(j+1)\,\#\{k\in J:k<j\}+\sum_{i\in K}(p-i-1)+\sum_{j\in K}(j+1)+h.
\]
Evaluating these sums in each chamber, with
\[
\begin{aligned}
A_0&=\tfrac{p(p+1)(p+2)}{6}, & B_0&=\tfrac{p(p+1)(p+5)}{6}, & C_0&=\tfrac{p(p+1)(2p+7)}{6},
& D_0&=\tfrac{p(p-1)(p+1)}{6},\\
E_0&=\tfrac{p(p-1)(2p+5)}{6}, & F_0&=\tfrac{(p-1)(p^2+4p+6)}{6}, & G_0&=\tfrac{(p-1)(p+1)(p+3)}{3},
\end{aligned}
\]
yields the complete orientation table (each orientation pair summing to $N$):
\begin{center}
\begin{tabular}{@{}lccccc@{}}\toprule
chamber & $N$ & $\#L(x\prec z)$ & $\#L(z\prec x)$ & $\#L(y\prec z)$ & $\#L(z\prec y)$\\ \midrule
right addition & $N_+$ & $C_0$ & $A_0$ & $B_0$ & $2A_0$\\
left addition & $N_+$ & $2A_0$ & $B_0$ & $A_0$ & $C_0$\\
right deletion & $N_-$ & $E_0$ & $F_0$ & $D_0$ & $G_0$\\
left deletion & $N_-$ & $G_0$ & $D_0$ & $F_0$ & $E_0$\\ \bottomrule
\end{tabular}
\end{center}
The displayed witnesses read off the table: $B_0/N_+=\frac{p+5}{3(p+3)}$ (balanced since
$p+5\ge p+3$ and $p+5\le 2p+6$); $2A_0/N_+=\frac{2(p+2)}{3(p+3)}$ (since
$2(p+2)\ge p+3$ and $2(p+2)\le2(p+3)$); $E_0/N_-=\frac{p(2p+5)}{3(p+1)(p+2)}$ (lower
bound $\iff p^2+2p-2\ge0$, upper $\iff p+4\ge0$); and
$F_0/N_-=\frac{p^2+4p+6}{3(p+1)(p+2)}$ (the same two inequalities exchanged). All hold
for $p\ge2$; reverse orientations are complements, hence balanced as well.
\end{proof}

This is the first chamber in which the triple-diagonal machinery operates without full
relabeling symmetry: the orbit argument no longer applies, and balance is carried by the
explicit fixed witnesses above.

\section{Closed-third avoidance: the rigidity theory}\label{sec:rigidity}
% Sources: closed-third-strong-order-rigidity + corrected-strong-order-theorem-under-
% closed-third-avoidance (r12); maximum-chain-heavy-gap-rigidity + heavy-gap-rigidity-
% under-closed-third-avoidance; heavy-gap-alignment-corollary-under-closed-third-
% avoidance + heavy-gap-probabilistic-alignment; maximum-chain-defect-windows-are-
% nonsingleton. All banked with independent re-derivations (ledger rows S4-S6).
From here on, the argument is a proof by contradiction, and this section builds the
contradiction hypothesis into a rigid structure. Say that $P$ satisfies
\emph{closed-third avoidance} if
\[
\pr(u\prec v)\notin[1/3,\,2/3]\qquad\text{for every orientation of every incomparable
pair }\{u,v\}.
\]
The Three-Defect Chain Theorem is exactly the statement that no
ordinal-sum-indecomposable non-chain poset with $|P\setminus C|\le3$ satisfies it.

\begin{theorem}[Strong order]\label{thm:strongorder}
Under closed-third avoidance, the relation $u\prec v\iff\pr(u\prec v)>2/3$ is a strict
total order on $P$ extending the poset order. If $P$ is not a chain, some adjacent pair
of this total order is incomparable in $P$.
\end{theorem}

\begin{proof}
Totality: for comparable pairs the orientation probability is $1$; for incomparable
pairs, $\pr(u\prec v)+\pr(v\prec u)=1$, and avoidance places exactly one of the two above
$2/3$. Transitivity: if $A=\{L:x<_Ly\}$ and $B=\{L:y<_Lz\}$ have $|A|,|B|>2N/3$, then
$|A\cap B|\ge|A|+|B|-N>N/3$, and every extension in $A\cap B$ has $x<_Lz$; so
$\pr(x\prec z)>1/3$. If $z<_Px$ this event would be empty, a contradiction, so either
$x<_Pz$ (whence $x\prec z$ directly) or $x\parallel z$, in which case avoidance excludes
$[1/3,2/3]$ and the established lower bound forces $\pr(x\prec z)>2/3$. For the adjacent
pair: if every adjacent pair of the total order were comparable in $P$, transitivity of
$P$ would make every pair comparable, so $P$ would be a chain.
\end{proof}

\begin{remark}[The closed interval is exact]
The open-interval analogue of avoidance is attainable: the four-element poset
$r<a<b$, $r<c$ has three extensions and its two incomparable pairs sit exactly at
$1/3$ and $2/3$. The conjecture's closed target $[1/3,2/3]$ is therefore the correct
boundary convention, and it is the closed hypothesis that the assembly refutes.
\end{remark}

For a defect $u\notin C$ and an extension $L$, let $G_u(L)\in\{0,\dots,m\}$ be the
\emph{gap} of $u$: the number of chain elements preceding $u$ in $L$. Legality confines
$G_u$ to the window $I_u=[a_u,b_u]$, and the chain-cut events are gap events:
$\{L:u<_Lc_t\}=\{G_u<t\}$. Write $\mu_u(r)=\#\{L:G_u(L)=r\}$.

\begin{lemma}[Nonsingleton windows]\label{lem:nonsingleton}
If $C$ is a maximum chain and $u\notin C$, then $|I_u|\ge2$.
\end{lemma}

\begin{proof}
If $I_u=\{r\}$, then $u$ is comparable with every chain element and inserts at one
position, so $C\cup\{u\}$ is a chain of size $m+1$, contradicting maximality.
\end{proof}

\begin{theorem}[Heavy-gap rigidity]\label{thm:heavygap}
Under closed-third avoidance, every defect $u$ has a unique gap $\rho_u$ with
$\mu_u(\rho_u)>N/3$, and both strict tails satisfy
$\sum_{r<\rho_u}\mu_u(r)<N/3$ and $\sum_{r>\rho_u}\mu_u(r)<N/3$.
\end{theorem}

\begin{proof}
The cumulative counts $F(t)=\#\{L:G_u<t\}$ run from $0$ to $N$ as $t$ crosses the window;
by Lemma~\ref{lem:nonsingleton} the window has interior thresholds, and each interior
value $F(t)$ is the orientation count of a genuine incomparable pair $\{u,c_t\}$, hence
avoids $[N/3,2N/3]$. Let $t$ be the last threshold with $F(t)<N/3$ (window endpoints
supply $F=0$). The next value cannot lie in the forbidden closed interval, so
$F(t+1)>2N/3$, and the single gap $\rho_u=t$ between them has
$\mu_u(\rho_u)=F(t+1)-F(t)>N/3$. Its lower tail is $F(t)<N/3$ and its upper tail is
$N-F(t+1)<N/3$; any other gap sits inside one of the tails, so no second gap reaches
$N/3$. Endpoint heavy gaps use the empty outside tail. % banked: maximum-chain-heavy-
% gap-rigidity; the open-third analogue is DISPROVED (2-chain + isolated point realizes
% CDF 0,1/3,2/3,1), recorded in maximum-chain-unique-heavy-gap-rigidity-result.
\end{proof}

\begin{theorem}[Heavy-gap alignment]\label{thm:alignment}
Let $R=P[P\setminus C]$ be the induced defect order and let $u,v$ be defects with
$\rho_u<\rho_v$. Under closed-third avoidance: (a) $v<_Ru$ is impossible; (b) if
$u\parallel_Rv$ then $\pr(u\prec v)>2/3$.
\end{theorem}

\begin{proof}
(a) If $v<_Ru$, isotonicity gives $G_v\le G_u$ in every extension, so
$\{G_v=\rho_v\}\subseteq\{G_u\ge\rho_v\}\subseteq\{G_u>\rho_u\}$, whence
$\mu_v(\rho_v)\le\sum_{r>\rho_u}\mu_u(r)<N/3$, contradicting heaviness.
(b) On the complement of $\{G_u>\rho_u\}\cup\{G_v<\rho_v\}$ one has
$G_u\le\rho_u<\rho_v\le G_v$, a strict gap inequality, which forces $u<_Lv$. Hence
$\pr(v\prec u)<1/3+1/3=2/3$; avoidance excludes $[1/3,2/3]$, so $\pr(v\prec u)<1/3$ and
$\pr(u\prec v)>2/3$.
\end{proof}

\section{The three-defect model and the case grid}\label{sec:model}
% Sources: uniform-gap-map bijection + v-lambda-uniform-three-defect-foundation-audit-
% round-28-corr (fiber model); three-defect-order-shape-classification; three-chain-
% defect-width-two-reduction + strict-reduction-atlas-for-three-chain-defects;
% canonical-ordinal-sum-factorization-theorem-for-finite-poset; round-29-two-equal-
% cell-partition (ledger rows S2, S7, S8).
Now let $|P\setminus C|=3$. A \emph{gap map} assigns to each defect $u$ a gap
$g(u)\in I_u$, isotone with respect to $R$; extensions decompose into fibers over
feasible gap maps, and with $B_r=g^{-1}(r)$ the fiber weight is
$W(g)=\prod_re\!\left(P[B_r]\right)$, so that $N=\sum_gW(g)$. The block factors here are
posets on at most three points, so $W\in\{1,2,3,6\}$; the weighted diagonals of
Sections~2 and~6 are the two- and three-defect instances of this formula.

\begin{proposition}[Shape classification]\label{prop:shapes}
Up to order duality, the induced defect order $R$ on three elements is one of: the
antichain $A$; the one-relation poset $Q$ ($x<y$, $z$ incomparable to both); the shape
$V$ ($p<q$, $p<r$, $q\parallel r$) or its dual $\Lambda$; or the three-chain $T$.
\end{proposition}

\begin{proposition}[Three-chain reduction]\label{prop:threechain}
If $R$ is the chain $x<y<z$, then $C\cup\{x<y<z\}$ is a two-chain cover of $P$, so $P$
has width at most two and the width-two theorem of Linial \cite{linial} (sharpened by Sah
\cite{sah}; the direct-sum exceptional family there is excluded by ordinal-sum
indecomposability) supplies a balanced pair.
\end{proposition}

By Theorem~\ref{thm:heavygap} each defect carries a unique heavy gap; three values
realize exactly one of three patterns: \emph{all-distinct} (D), \emph{exactly two equal}
(2E), or \emph{all equal} (CH). Together with Proposition~\ref{prop:shapes} this yields
the twelve-cell grid
\[
\{A,\,V,\,\Lambda,\,Q\}\times\{\mathrm{D},\,\mathrm{2E},\,\mathrm{CH}\},
\]
and the Three-Defect Chain Theorem is proved by showing every cell is empty under
closed-third avoidance. Sections~10--12 do this shape by shape. We will also use:

\begin{theorem}[Canonical factorization; balance under ordinal sums]\label{thm:factor}
Every finite nonempty poset factors uniquely as the ordinal sum of the induced subposets
on the connected components of its incomparability graph, each factor
ordinal-sum-indecomposable; and a uniformly random extension of the sum restricts to
independent uniform extensions of the factors, so every incomparable pair has the same
orientation probabilities in the sum as in its factor.
\end{theorem}

\section{The antichain shape}\label{sec:antichain}
% Sources: antichain-defect-distinct-heavy-index-exclusion (+ corrected derivation and
% debt resolution, 2D); centered-three-interval-ch-chamber-specific-global-decision +
% deepcheck-repaired-ch-global-residuals-next + exceptional-weak-chamber-ch-certificate
% (rows A-D, A-2E, A-CH).
\begin{theorem}[Heavy-gap coincidence for antichain defects]\label{thm:antichaincoincide}
Let $x,y,z$ be pairwise incomparable defects. Under closed-third avoidance,
$\rho_x=\rho_y=\rho_z$.
\end{theorem}

\begin{proof}[Proof sketch from the banked derivation]
With no relations among the defects, the feasible gap maps form the product
$I_x\times I_y\times I_z$, and the exact tied-fiber weights make each defect's marginal
$\mu_u$ largest precisely on the common intersection $K=I_x\cap I_y\cap I_z$: a common
gap acquires the tied-block multiplicities, and moving off $K$ strictly sheds them.
Pairwise overlap of the three windows (Lemma~\ref{lem:overlap} applied pairwise) gives
$K\neq\emptyset$ by Helly's theorem in one dimension. Heavy-gap uniqueness
(Theorem~\ref{thm:heavygap}) then pins all three heavy gaps to the same point of $K$.
\end{proof}

This empties A-D and A-2E simultaneously --- any non-coincident pattern is impossible.
The surviving cell is the common-heavy antichain, which is destroyed by an exact
inequality. Translate the common heavy gap to $0$, write $I_u=[-a_u,b_u]$, and let
$T_{\sigma}$ count extensions by the relative order $\sigma$ of $(x,y,z)$ (indices
$1,2,3$), so $N=\sum_\sigma T_\sigma$. Under avoidance the strong order
(Theorem~\ref{thm:strongorder}) orients the three defect pairs with probabilities
$>2/3$; labeling so that $x\prec y\prec z$, adding the three inequalities
$\#(x<y),\#(y<z),\#(x<z)>2N/3$ and subtracting $2N$ gives
\[
Q_A\;:=\;T_{123}-T_{231}-T_{312}-2T_{321}\;>\;0 .
\]
Meanwhile $L_1$ (the lower tail of the lowest window's defect) and $H_3$ (the upper tail
of the highest) are strict heavy-gap tails, so $3L_1<N$ and $3H_3<N$.

\begin{theorem}[Centered three-interval CH inequality]\label{thm:CH}
For every nonnegative integer endpoint vector with $I_u=[-a_u,b_u]$: if $3L_1<N$ and
$3H_3<N$, then $Q_A\le0$.
\end{theorem}

\begin{proof}[Proof by certificate corpus]
Set $S_L=N-3L_1$, $S_H=N-3H_3$. The weak orderings of the two endpoint triples
$(a_1,a_2,a_3)$ and $(b_1,b_2,b_3)$ partition parameter space into $13\times13=169$
chambers, on each of which $N$, $L_1$, $H_3$, $Q_A$ are polynomials in the chamber's
nonnegative coordinates. For $168$ chambers there are exact rational $\alpha,\beta\ge0$
with
\[
-Q_A\;=\;\alpha S_L+\beta S_H+\mathrm{Rem},
\]
where every coefficient of $\mathrm{Rem}$ is nonnegative --- so $S_L,S_H>0$ forces
$Q_A\le0$. The single exceptional chamber ($a$-pattern $12|0$, $b$-pattern $01|2$)
carries the closed-form parametrization $a_2=a_3=p$, $a_1=p+1+u$, $b_1=b_2=q$,
$b_3=q+1+v$, in which
\[
Q_A=-\tfrac{p+q+1}{2}\,R,\qquad
R=p^2+2pq-pu-pv+3p+q^2-qu-qv+3q-2uv-4u-4v,
\]
and the exact expansion of $S_L+S_H$ is at most $-2$ when $u+v\ge1$ and strictly
negative when $u=v=0$, $p+q\ge1$. Hence $S_L>0$, $S_H>0$ forces $p=q=u=v=0$, where
$Q_A=0$. The full certificate table and its exact-coefficient audit are proof
object~(i) of Section~\ref{sec:proofobjects}.
\end{proof}

Together, $Q_A>0$ and $Q_A\le0$ are the contradiction: the cell A-CH is empty.

\section{The $V$ and $\Lambda$ shapes}\label{sec:vlambda}
% Sources: v-lambda-uniform-three-defect-foundation-audit-round-28-corr (V-D, L-D);
% every-realizable-gated-common-heavy-v-configuration-fails-cl + v-common-heavy-vacuity
% + lambda-common-heavy-vacuity-by-duality (V-CH, L-CH); round-29-two-equal-cell-
% partition; full-cone-te-2-1-theorem-closing-v-2e-tied-siblings + gated-te21-chamber-
% atlas-and-exact-certificates + exact-d-f-facet-certificate-atlas-for-guarded-te-2-1 +
% normalized-d-f-0-te-2-1-facet-certificate-atlas + independent-full-cone-te-2-1-
% certificate-replay (V-2E); lambda-l-2e-exact-reflection-transfer(+-from-full-cone) +
% independent-lambda-reflection-replay (L-2E).
Let the defect order be $V$: $p<q$, $p<r$, $q\parallel r$, with windows
$I_p=[A,B]$, $I_q=[C_0,D]$, $I_r=[E,F]$ (siblings exchanged so $D\le F$; realizability
gives $A\le C_0$, $A\le E$, $B\le D$, $E\le D$). The feasible states and
weights are
\[
\Gamma=\{(i,j,k)\in I_p\times I_q\times I_r:\ i\le j,\ i\le k\},\qquad
W(i,j,k)=1+\mathbf1_{\{j=k\}} .
\]

\begin{theorem}[Sibling coincidence; V-D empty]\label{thm:vd}
Under closed-third avoidance, the sibling heavy gaps coincide: $\rho_q=\rho_r=D$.
Consequently the all-distinct pattern is impossible for $V$.
\end{theorem}

\begin{proof}
The sibling marginal $M_q(j)=\sum_{i\in I_p,i\le j}|I_r\cap[i,m]|
+\mathbf1_{\{j\in I_r\}}|I_p\cap(-\infty,j]|$ is nondecreasing on $I_q$, so its unique
heavy gap is the right endpoint $D$. The marginal $M_r$ is nondecreasing through $D$;
beyond $D$ its first term is constant and its diagonal term vanishes, while at $D$ the
diagonal contributes $|I_p|>0$, so $M_r(D)>M_r(k)$ for all $k>D$. Uniqueness pins
$\rho_r=D=\rho_q$.
\end{proof}

\begin{theorem}[V-CH is vacuous]\label{thm:vch}
No realizable $V$-configuration has a common heavy gap.
\end{theorem}

\begin{proof}
For $i\in I_p$ put $Q_i=|I_q\cap[i,m]|$, $R_i=|I_r\cap[i,m]|$, $T_i=|I_q\cap I_r\cap[i,m]|$;
fixing $i$, the feasible $(j,k)$ are the two suffixes and the diagonal contributes one
extra extension, so the marginal mass of $p$ is $M_p(i)=Q_iR_i+T_i$ --- a product of
suffix counts, hence \emph{nonincreasing} in $i$. The marginal of $q$,
$M_q(j)=\sum_{i\in I_p,\,i\le j}|I_r\cap[i,m]|+\mathbf1_{\{j\in I_r\}}|I_p\cap(-\infty,j]|$,
is \emph{nondecreasing} on $I_q$ (both terms are; the diagonal indicator switches on at
$E$ and stays on through $D\le F$). Now suppose $h$ is heavy for all three defects. If
$h<D$, then the upper tail of $q$ at $h$ contains $M_q(D)\ge M_q(h)>N/3$, contradicting
the tail bound of Theorem~\ref{thm:heavygap}; so $h=D$. If $h>A$, the lower tail of $p$
at $h$ contains $M_p(A)\ge M_p(h)>N/3$; so $h=A$. Hence $A=D$, and realizability's
$A\le C_0\le D$ forces $C_0=D$: the window $I_q=\{D\}$ is a singleton, contradicting
Lemma~\ref{lem:nonsingleton}.
\end{proof}

\begin{theorem}[Two-equal partition]\label{thm:partition}
The exactly-two-equal dispatch for the three shapes with relations partitions into $14$
cells: $5$ for $V$, $5$ for $\Lambda$, $4$ for $Q$. Eight are closed by the coincidence
and alignment exclusions above and their duals; the survivors are the \emph{tied-siblings}
$V$ cell ($\rho_q=\rho_r>\rho_p$), its $\Lambda$ dual ($\rho_q=\rho_r<\rho_p$), and four
$Q$ cells treated in Section~\ref{sec:q}.
\end{theorem}

In the tied-siblings cell the defect $p$'s marginal decreases in $i$ (larger $i$ shrinks
the feasible $(j,k)$-region), so $\rho_p=A$, the left endpoint of $I_p$. The two pinned
quantities are the strict heavy-gap tails
\[
H=\#\{\text{weighted states with }i>A\},\qquad
L=\#\{\text{weighted states with }j<D\},
\]
and avoidance gives $3H<N$ and $3L<N$. The cell is emptied by:

\begin{theorem}[Full-Cone TE$(2{=}1)$ Theorem]\label{thm:fullcone}
Every canonical realizable tied-siblings instance satisfies $N\le3\max(H,L)$.
\end{theorem}

\begin{proof}[Proof by certificate corpus]
Canonical sibling exchange gives $D\le F$, splitting the cone into the disjoint facets
$D<F$ and $D=F$. Translating $D$ to $0$, realizability confines the endpoint parameters
to a four-variable integer cone whose weak orderings give exactly $31$ chambers on each
facet (of the $75$ ordered weak partitions of four labels, $31$ satisfy $x>y$). On the
strict facet, each chamber carries two exact signed certificates ($62$ in total)
expressing the branch identities
$3H-N=\alpha(N-3L)+\beta(N-3H)+R_H$ and $3L-N=\alpha(N-3L)+\beta(N-3H)+R_L$ with
$\alpha,\beta\ge0$ rational and residuals coefficientwise nonnegative, which force
$N\le3\max(H,L)$ chamber by chamber. On the boundary facet $D=F$ the stronger inequality
holds: all $31$ residual identities $3L-N=R_L$ have coefficientwise nonnegative $R_L$
($391$ nonzero monomials, minimum nonzero coefficient $1/2$), so $N\le3L$ outright. The
full-cone partition of the realizable domain into the two facets is exhaustive and
disjoint. Certificates, generators, and independent replays are proof objects~(ii)--(iv)
of Section~\ref{sec:proofobjects}.
\end{proof}

Since $\max(H,L)\ge N/3$ contradicts both strict tails, V-2E is empty --- and with it
the last $V$ cell. The $\Lambda$ shape now falls by an exact duality.

\begin{theorem}[Reflection transfer; $\Lambda$ empty]\label{thm:reflection}
Reflection of chain gaps by $R_m(x)=m-x$ (reversing the chain) bijects the feasible
$\Lambda$ states $\{(i,j,k):j\le i,\ k\le i\}$ with the canonical $V$ states, preserves
the fiber weight $1+\mathbf1_{\{j=k\}}$ and hence $N$, maps windows $[a,b]$ to
$[m-b,m-a]$, exchanges strict lower and upper tails, and sends every same-label
orientation numerator $S$ to $N-S$ --- so closed-third avoidance, realizability, and
heavy-gap data transfer exactly. Consequently: $\Lambda$-D is empty by
Theorem~\ref{thm:vd}, $\Lambda$-CH by Theorem~\ref{thm:vch}, and the tied-siblings
$\Lambda$ cell maps onto the $V$ TE$(2{=}1)$ cone, where Theorem~\ref{thm:fullcone}
contradicts the two transferred strict guards. All of $\Lambda$ is empty.
\end{theorem}

\section{The $Q$ shape}\label{sec:q}
% Sources: pairwise-distinct-heavy-q-exclusion + q-distinct-heavy-chamber-certificate;
% common-heavy-q-exclusion + q-common-heavy-closure-decision-round-23-repaired;
% q1-q2-two-equal-structural-exclusion (r30); q3-l-and-q4-r-diagonal-correction +
% corrected-q3-l-q4-r-full-diagonal-transfer + aligned-corrected-q3-l-q4-r-formulas +
% corrected-ordered-partition-audit-and-final-verdict-for-q-2e + independent-q-2e-34-
% certificate-replay (rows Q-D, Q-CH, Q-2E).
Let the defect order be $Q$: $x<y$, $z$ incomparable to both. The feasible states are
$(i,j,k)\in I_x\times I_y\times I_z$ with $i\le j$, weighted
$W=1+\mathbf1_{\{i=k\}}+\mathbf1_{\{j=k\}}$ (the full diagonal $i=j=k$ has weight $3$).

\begin{theorem}[Q-D empty]\label{thm:qd}
Under closed-third avoidance, pairwise-distinct heavy gaps are impossible: the three
interval chambers $b<c$, $b=c$, $c<b$ of the comparable pair's windows $[a,b],[c,d]$
exhaust the legal systems, and each either forces coincidence or degeneracy of heavy
indices or yields a genuine defect--chain cut whose orientation numerator lies strictly
inside $(N/3,\,2N/3)$.
\end{theorem}

\begin{theorem}[Q-CH empty]\label{thm:qch}
If $x,y,z$ share a heavy gap $r$, then one of the strict tails $L_x(r)$, $H_y(r)$ is the
orientation numerator of a genuine incomparable defect--chain pair and satisfies
$N/3<T<2N/3$; the complementary side's heavy atom ($>N/3$) supplies the upper bound.
\end{theorem}

\begin{theorem}[Q1, Q2 empty]\label{thm:q12}
In the two two-equal cells whose tied pair is the comparable pair $\{x,y\}$, the equal
heavy gaps $\rho_x=\rho_y$ are excluded structurally: the $Q$ state set has no
constraint $i\le k$, and the exact unique-mode separation of the $x$ and $y$ marginals
forces $\rho_x\neq\rho_y$.
\end{theorem}

The remaining cells Q3, Q4 tie $z$ with one of the comparable defects: Q3 is
$\rho_x=\rho_z<\rho_y$ (tied $\{x,z\}$, $y$'s heavy gap above), Q4 is
$\rho_x<\rho_y=\rho_z$. Write $X=[a,b]$, $Y=[c,d]$, $Z=[e,f]$ for the windows, with
lengths $n_x,n_y,n_z\ge2$ (Lemma~\ref{lem:nonsingleton}). The comparable pair's windows
meet in exactly one of three chambers, and the exact marginal mode analysis of the
banked classification disposes of each: if $b=c$ (one-point overlap), a nonsingleton
$Z$ creates a two-point maximal plateau in the tied marginals, so no \emph{unique}
heavy gap exists and the cell is infeasible; if $b<c$ (\emph{disjoint}), the cell is
killed by Lemma~\ref{lem:qdisjoint} below; if $c<b$ (\emph{genuine overlap}), the tied
coincidence forces the endpoint chamber ($Z=[e,c]$, $X=Y=[c,b]$ for Q3; $Z=[b,f]$,
$X=Y=[c,b]$ for Q4), which transfers to the guarded-$V$ cone of
Section~\ref{sec:vlambda} with a one-unit correction (Theorems~\ref{thm:qtransfer}
and~\ref{thm:q34}).

\begin{lemma}[Disjoint-branch closure]\label{lem:qdisjoint}
In the disjoint chamber $b<c$, both cells are impossible under closed-third avoidance.
\end{lemma}

\begin{proof}
Here mode uniqueness forces $Z=[b,c]$ with $\rho_x=b$, $\rho_y=c$, and the tie
determines the cell: $\rho_z=b$ (cell Q3) exactly when $n_y>n_x$, and $\rho_z=c$ (cell
Q4) exactly when $n_x>n_y$. The exact extension count is $N=n_zn_xn_y+n_x+n_y$. For Q3
select the genuine chain pair $\{x,c_b\}$ (eligible since $a<b$); its orientation
numerator is the strict lower tail of $x$ at its heavy gap,
\[
A=(n_x-1)(n_zn_y+1),\qquad
3A-N=n_zn_y(2n_x-3)+2n_x-n_y-3,
\]
which is positive for $n_y>n_x\ge2$: at $n_x=2$ it equals $(n_z-1)n_y+1>0$, and for
$n_x\ge3$ it is at least $n_y(4n_x-7)+2n_x-3>0$. The complementary mass contains the
heavy atom, $N-A\ge\mu_x(b)>N/3$, so $A<2N/3$. Thus $N/3<A<2N/3$, contradicting
avoidance. For Q4, symmetrically, the pair $\{c_{c+1},y\}$ (eligible since $c<d$) has
numerator $A=(n_y-1)(n_zn_x+1)$, the strict upper tail of $y$ at $c$, with
$3A-N=n_zn_x(2n_y-3)+2n_y-n_x-3>0$ for $n_x>n_y\ge2$ by the same two-case check, and
$A<2N/3$ from the heavy atom again.
\end{proof}

\begin{theorem}[Corrected full-diagonal transfer]\label{thm:qtransfer}
For every realizable corrected Q3-L or Q4-R overlap instance, gap reversal (Q3-L) or
direct relabeling (Q4-R) identifies the $Q$ state set with a canonical guarded-$V$ state
set, and the fiber weights agree except at the full diagonal, where $Q$ carries weight
$3$ against $V$'s $2$. The exact count relations are
\[
N_Q=N_V+1,\qquad H_Q=H_V+1,\qquad L_Q=L_V,
\]
with the extra unit on a single named sibling orientation.
\end{theorem}

\begin{theorem}[Q3, Q4 empty]\label{thm:q34}
Under closed-third avoidance the corrected overlap branches are impossible.
\end{theorem}

\begin{proof}[Proof from the certificate atlases]
Failure of both corrected bounds means $N_Q>3H_Q$ and $N_Q>3L_Q$; by
Theorem~\ref{thm:qtransfer} and integrality this is equivalent to
\[
N_V-3H_V\ge3\qquad\text{and}\qquad N_V-3L_V\ge0 .
\]
If $H_V\ge L_V$, Theorem~\ref{thm:fullcone} gives $3H_V\ge N_V$, contradicting the first
inequality. If $L_V>H_V$, the theorem gives $3L_V\ge N_V$ and only equality
$N_V=3L_V$ survives; the certificates exclude it. On strict-facet chambers whose residual
$R_L$ has a positive constant term, equality is immediate; on the zero-constant chambers
the signed identity $3L_V-N_V=\alpha(N_V-3L_V)+\beta(N_V-3H_V)+R_L$ with $\beta>0$
evaluates, at $N_V=3L_V$ and $L_V>H_V$, to a strictly positive right side against a zero
left side. On the boundary facet the equality-face audit gives the strict floor
$3L_V-N_V\ge1$ on every realizable point. Integrality then supplies the missing unit in
each case, including the tied case $L_V=H_V+1$, where $H_Q=L_Q=L_V$ and
$N_Q=N_V+1\le3L_V=3\max(H_Q,L_Q)$. Finally the \emph{aligned} family
($A=C_0$, $B=D=E$, $D<F$) admits closed forms with $n=D-A+1\ge2$, $m=F-D+1\ge2$:
\[
N_Q=\tfrac{mn(n+1)}2+n+1,\quad H_Q=\tfrac{mn(n-1)}2+n,\quad L_Q=\tfrac{mn(n-1)}2,
\]
whence $3H_Q-N_Q=mn(n-2)+2n-1\ge3$ directly. The corrected chamber-by-chamber and
boundary certificates, with their independent replay, are proof object~(v) of
Section~\ref{sec:proofobjects}.
\end{proof}

\section{The Three-Defect Chain Theorem}\label{sec:assembly}
% Source: three-defect-chain-theorem + three-defect-chain-theorem-audited-final-assembly
% (r38); the slug-cited acyclic dependency walk over the validated ASSEMBLY_LEDGER.
\begin{theorem}[Three-Defect Chain Theorem]\label{thm:threedefect}
Every finite ordinal-sum-indecomposable non-chain poset $P$ with a maximum chain $C$ and
$|P\setminus C|\le3$ has an incomparable pair $\{u,v\}$ with $1/3\le\pr(u\prec v)\le2/3$.
\end{theorem}

\begin{proof}
Assume closed-third avoidance. If $|P\setminus C|\le2$, Theorem~\ref{thm:twodefect}
already yields a balanced pair (defect $0$ contradicts non-chain; defect $1$ has width
two). Let $|P\setminus C|=3$. If the defects form a chain,
Proposition~\ref{prop:threechain} applies. Otherwise, by
Proposition~\ref{prop:shapes} the defect order is (up to duality) $A$, $V$, $\Lambda$,
or $Q$, and by Theorem~\ref{thm:heavygap} the heavy pattern is D, 2E, or CH. Every cell
of the resulting grid is empty:
A-D and A-2E by Theorem~\ref{thm:antichaincoincide}, A-CH by Theorem~\ref{thm:CH};
V-D by Theorem~\ref{thm:vd}, V-CH by Theorem~\ref{thm:vch}, V-2E by
Theorems~\ref{thm:partition} and~\ref{thm:fullcone}; all of $\Lambda$ by
Theorem~\ref{thm:reflection}; Q-D by Theorem~\ref{thm:qd}, Q-CH by
Theorem~\ref{thm:qch}, Q-2E by Theorem~\ref{thm:q12} (cells Q1, Q2),
Lemma~\ref{lem:qdisjoint} (the disjoint and one-point chambers of Q3, Q4), and
Theorem~\ref{thm:q34} (their overlap chambers). The
contradiction hypothesis is untenable, so some incomparable pair lands in the closed
middle third. The dependency walk is acyclic: each cell closure consumes only the
rigidity theory of Section~\ref{sec:rigidity}, the model of Section~\ref{sec:model}, and
its own certificates.
\end{proof}

\begin{corollary}[Defect at most three]\label{cor:defect3}
Every finite non-chain poset of defect at most three has an incomparable pair
$\{u,v\}$ with $1/3\le\pr(u\prec v)\le2/3$.
\end{corollary}

\begin{proof}
Chains of the canonical factorization $P=P_1\oplus\cdots\oplus P_k$
(Theorem~\ref{thm:factor}) concatenate: a maximum chain of $P$ is the concatenation of
maximum chains of the factors, so the defect of $P$ is the sum of the defects of the
factors. Since $P$ is not a chain, some indecomposable factor $P_i$ is not a chain, and
its defect is between $1$ and $3$; Theorem~\ref{thm:threedefect} gives a balanced pair
in $P_i$, and Theorem~\ref{thm:factor} preserves its orientation probabilities in $P$.
\end{proof}

\section{Proof objects and verification}\label{sec:proofobjects}
% The certificate-corpus discipline: exact rational symbolic identities, coefficientwise
% checks, independent replays, SHA-256 anchors. Corpus hashes verbatim from the banked
% artifacts (r26-r38).
The chamber theorems above are proved by finite families of \emph{exact symbolic
identities}: on each chamber, the quantities $N,H,L$ (or $N,L_1,H_3,Q_A$) are
polynomials with rational coefficients in the chamber's nonnegative integer coordinates,
and each certificate is an identity between such polynomials whose validity is checked
coefficientwise in exact rational arithmetic --- no floating point, no sampling, no
interpolation in the checker. The corpora are machine-readable artifacts; each was
re-derived by an \emph{independent replay} (a separately written checker that
reconstructs the identities from the raw definitions), and the artifacts are anchored by
SHA-256 hashes. The load-bearing corpora are:
\begin{itemize}
\item[(i)] \textbf{CH atlas} (Theorem~\ref{thm:CH}): $169$ weak-order chambers, $168$
exact $(\alpha,\beta)$ certificates plus the exceptional-chamber identity; exact
coefficient audit of every residual. Decision artifact:\\
{\footnotesize\texttt{9eb6031fecf65abb51f16ed73adc8d166bb3582ecd94467f4415220f8d94a87e}}
\item[(ii)] \textbf{TE$(2{=}1)$ strict-facet atlas} (Theorem~\ref{thm:fullcone}): $31$
chambers, $62$ signed certificates; replay executed $49{,}152$ exact identity-component
checks and $142{,}125$ numeric spot checks. Atlas / exact coefficient data:\\
{\footnotesize\texttt{3541901a46c6e9603933d95546d3adb075df15570e6cc6134c7a67ed2c8777a7}}\\
{\footnotesize\texttt{59e2b841ce2cb653d900a73092165f1b78547f4bf6023c4b04adf677a878174d}}
\item[(iii)] \textbf{TE$(2{=}1)$ boundary-facet atlas} ($D{=}F$): $31$ chambers, $391$
nonnegative residual monomials, equality-face audit with minimum residual $1$ on all
$\{0,1\}$-faces and $2{,}304$ finite-envelope rows. Atlas / exact coefficient data:\\
{\footnotesize\texttt{e5cdb27fb431b3e0ef84c6edbc95dd2e525585ec665278cf67564aab6ab3f219}}\\
{\footnotesize\texttt{7aad8e7378dfe5475ba1dc3b8283e8250adf64dafe3359a0a0460b23c6288a5a}}
\item[(iv)] \textbf{Full-cone replay and $\Lambda$ reflection replay}
(Theorems~\ref{thm:fullcone},~\ref{thm:reflection}): fresh checkers partitioned $6{,}048$
canonical rows exhaustively into $5{,}292$ strict and $756$ boundary rows with no gap or
overlap, and verified the reflection on $173{,}613$ interval triples, $6{,}355{,}899$
tail identities, and $3{,}751{,}704$ chain-cut identities, plus actual linear-extension
duality on all realizable cases through $m=3$. Replay record / reflection log:\\
{\footnotesize\texttt{6fce15f9b70e629b175f0175a07fe76f69fabc7b0c118ee431491f219271bca2}}\\
{\footnotesize\texttt{a9c266e2bd909c33368a10602d4588ea874c8d6c5c02de44ad53633feef0e193}}
\item[(v)] \textbf{Corrected $Q$ replay} (Theorem~\ref{thm:q34}): $31$ strict chambers,
$62$ signed certificates, $31$ boundary identities, $7{,}776$ strict and $2{,}304$
boundary coverage rows, no negative coefficient and no failed identity. Replay record:\\
{\footnotesize\texttt{58c6d5909596794f2561882c625afb3193d077b92f9918dc0fce656a0cbd1e8e}}
\end{itemize}
All hashes are SHA-256 of the artifact files exactly as published. The complete corpora
--- certificate tables, exact rational coefficient data, generators, independent
checkers, and a \texttt{SHA256SUMS} manifest covering every file --- accompany this
preprint in the public artifact repository
\texttt{github.com/emios-labs/emi-papers} (directory \texttt{defect-chains}, tagged
release matching this draft), together with a supplementary document
assembling the full banked derivations behind the statement-level theorems
(Theorems~\ref{thm:antichaincoincide},~\ref{thm:partition},~\ref{thm:qd},~\ref{thm:qch},
\ref{thm:q12}, and~\ref{thm:qtransfer}). The replays corroborate and gate the symbolic
corpora; the theorems rest on the exact identities themselves, which a referee can
re-expand from the published artifacts with any computer algebra system.

\section{Discussion: the defect ladder ahead}\label{sec:discussion}
% Updated 2026-08-18: the r10 obstruction paragraph is retained as motivation --- those
% blocked reductions are why the direct cell-closure architecture was necessary.
The reduction routes one might hope for are provably blocked at three defects: no
re-presentation as a two-defect chain exists inside a connected three-defect core, the
canonical ordinal-sum factorization produces no usable factor there, and deletion fails
quantitatively (restriction of the uniform distribution to $P-u$ is not uniform after
reinsertion). The direct architecture of Sections~8--13 --- avoidance rigidity, an
exhaustive case grid, and exact certificate corpora for the surviving cells --- was
forced by these obstructions, and it is the contribution we expect to scale.

Three of its pillars are already defect-count-free: the strong-order theorem, heavy-gap
rigidity, and the Helly coincidence mechanism of Theorem~\ref{thm:antichaincoincide}
apply verbatim to any number of pairwise-incomparable defects, and canonical
factorization reduces the general conjecture to indecomposable posets. What is not yet
defect-count-free is the endgame: at $k$ defects the shape classification grows, and the
common-heavy antichain cell needs a $k$-free replacement for
Theorem~\ref{thm:CH}. The concrete open target is a \emph{weighted-gap split-tail
inequality}: for antichain defects $x,y$ sharing the heavy gap under closed-third
avoidance, bound $|\pr(G_x<G_y)-\pr(G_x>G_y)|$ from the joint weighted fiber structure.
Marginals alone provably do not suffice --- independent marginals with below-tail
$1/3$/heavy $2/3$ against heavy $2/3$/above-tail $1/3$ reach $\approx0.78$ --- so any
proof must use the tied-diagonal weights and the coupling through the remaining defects.
A bound of $1/3$ for some pair would collapse the entire arbitrary-$k$ antichain slice
at once.

Other open questions: the defect-four ladder (new shapes, including the first
width-four cores); whether any aligned member of Section~6 attains balance constant
exactly $1/3$ (cf.\ Theorem~\ref{thm:rigidity}); and the interval-order connection ---
incomparable-defect $k$-defect chains are interval orders with $k$ nontrivial intervals,
so the antichain slice is a first approximation to the open interval-order case.

\begin{thebibliography}{99}
\bibitem{kislitsyn} S.~S. Kislitsyn, \emph{Finite partially ordered sets and their
associated sets of permutations}, Matematicheskie Zametki \textbf{4} (1968), 511--518.
\bibitem{fredman} M.~Fredman, \emph{How good is the information theory bound in sorting?},
Theoretical Computer Science \textbf{1} (1976), 355--361.
\bibitem{linial} N.~Linial, \emph{The information-theoretic bound is good for merging},
SIAM J. Computing \textbf{13} (1984), 795--801.
\bibitem{kahn-saks} J.~Kahn and M.~Saks, \emph{Balancing poset extensions}, Order
\textbf{1} (1984), 113--126.
\bibitem{aigner} M.~Aigner, \emph{A note on merging}, Order \textbf{2} (1985), 257--264.
\bibitem{ghp} B.~Ganter, G.~H\"afner, and W.~Poguntke, \emph{On linear extensions of
ordered sets with a symmetry}, Discrete Mathematics \textbf{63} (1987), 153--156.
% VERIFIED 2026-08-18 (web, produced source): GHP = Discrete Math 63 (1987) 153-156,
% cycle-free + nontrivial automorphism -> 1/2 pair; Peczarski 2006 = GPC (hence 1/3-2/3)
% for any poset with a nontrivial automorphism. Both cited accordingly in the intro.
\bibitem{brightwell-semi} G.~Brightwell, \emph{Semiorders and the 1/3--2/3 conjecture},
Order \textbf{5} (1989), 369--380.
\bibitem{tgf} W.~T. Trotter, W.~G. Gehrlein, and P.~C. Fishburn, \emph{Balance theorems
for height-2 posets}, Order \textbf{9} (1992), 43--53.
\bibitem{bft} G.~Brightwell, S.~Felsner, and W.~T. Trotter, \emph{Balancing pairs and the
cross product conjecture}, Order \textbf{12} (1995), 327--349.
\bibitem{peczarski-thin} M.~Peczarski, \emph{The Gold Partition Conjecture for 6-thin
posets}, Order \textbf{25} (2008), 91--103.
\bibitem{peczarski-aut} M.~Peczarski, \emph{The gold partition conjecture}, Order
\textbf{23} (2006), 89--95. % TODO(verify): exact home of the automorphism theorem
\bibitem{deloof} K.~De Loof, B.~De Baets, and H.~De Meyer, \emph{Counting linear
extension majority cycles in partially ordered sets on up to 13 elements}, Computers \&
Mathematics with Applications \textbf{59} (2010), 1541--1547.
\bibitem{zaguia-nfree} I.~Zaguia, \emph{The 1/3--2/3 conjecture for $N$-free ordered
sets}, Electronic J. Combinatorics \textbf{19}(2) (2012), P29.
\bibitem{zaguia-polytree} I.~Zaguia, \emph{The 1/3--2/3 conjecture for ordered sets whose
cover graph is a forest}, Order \textbf{36} (2019), 335--347.
\bibitem{olson-sagan} E.~J. Olson and B.~E. Sagan, \emph{On the 1/3--2/3 conjecture},
Order \textbf{35} (2018), 581--596.
\bibitem{sah} A.~Sah, \emph{Improving the $\frac13$--$\frac23$ conjecture for width two
posets}, Combinatorica \textbf{41} (2021), 99--126.
\bibitem{dolores-cuenca} E.~Dolores-Cuenca, A.~Guzm\'an-S\'aenz, and S.~Kim, \emph{The
gold partition conjecture and the lexicographic sum of posets}, arXiv:2410.12494.
\bibitem{chan-pak} S.~H. Chan and I.~Pak, survey. % TODO(verify): exact reference
\bibitem{gupta} R.~Gupta, \emph{Balance constants, majority cycles, and the Gold
Partition Conjecture through fourteen elements}, arXiv:2607.23926 (2026).
\end{thebibliography}

\end{document}
